% Front matter: edition notice (PUB-REF-0001 AN-20) and About This Book
% (AN-21 F-LOT: Who Should Read This Book / Conventions / Organization /
% Other Initech 123 Documentation / After You Complete This Book; annex
% Sec 5).  All boilerplate newly written (VO-12).
% GATES: test-pubs-build
\Gates{test-pubs-build}
\begin{TitleVerso}
\IsThisGoodForTheCompany\par
{\imLegalFont
\textbf{The following paragraph does not apply to any country where such
provisions are inconsistent with local law:} INITECH SYSTEMS CORPORATION
PROVIDES THIS PUBLICATION ``AS IS'' WITHOUT WARRANTY OF ANY KIND, EITHER
EXPRESS OR IMPLIED, INCLUDING, BUT NOT LIMITED TO, THE IMPLIED WARRANTIES OF
MERCHANTABILITY AND FITNESS FOR A PARTICULAR PURPOSE.  Some states do not
allow the exclusion of implied warranties, so the above exclusion may not
apply to you.\par
This publication may contain technical inaccuracies or typographical errors.
Its contents are revised from time to time, and the revisions appear in new
editions.  Initech Systems Corporation may change the products and programs
described here at any time without notice.\par
Copies of Initech publications, and technical information about Initech
products, are available from your authorized Initech dealer or your Initech
representative.\par}

{\footnotesize
Comments on this book are welcome; address them to: \InitechAddressLine.
\draft{PUB-REF-0001 O-5 and D-12: the address is to be checked against real
businesses before print}
Initech may use or distribute the information you supply in any way it
believes appropriate without incurring any obligation to you.\par
The software described in this publication is furnished under the Initech
Program License Agreement and may be used or copied only as that agreement
permits.\par}
\end{TitleVerso}

\FrontChapter{About This Book}
This book, the \doctitle{Initech 123 Reference}, describes every command,
every @function and every key of Initech 123, the worksheet program of
InitechOS.  Use it when you want the complete facts about a command: what
it does, what it asks you for, and what it leaves in the worksheet.

\section{Who Should Read This Book}
This book is for anyone who uses Initech 123.  It assumes that you know how
to start Initech 123 and move around a worksheet, as the
\doctitle{Setting Up Initech 123, Tutorial, Quick Start and Sample
Applications} volume shows.  Read Chapter~1 first; after that, read the
chapters in any order you like.

\section{Conventions}
This book uses the following conventions.
\begin{itemize}
\item Names of keys are in small capitals.  A function key that has a name
  in Initech 123 is given by its name followed by the key in parentheses.\\
  Example: Press \fkey{Calc}{F9} to recalculate the worksheet.
\item Two key names joined by a hyphen mean: hold down the first key and
  press the second.\\
  Example: Press \kname{CTRL-BREAK} to return to \textsf{READY} mode.
\item Two key names separated by a space mean: press the first key,
  release it, then press the second.\\
  Example: Press \kname{SLASH} \key{R} to choose the Range commands.
\item Text that you type is in this typeface.\\
  Example: Type \type{@SUM(A1..A4)} and press \kname{ENTER}.
\item Commands are named by the menu words you select, after a slash, as
  they appear in the control panel.\\
  Example: \slashcmd{/Range Format Currency}
\item @function arguments are in italics.\\
  Example: \fn{@ROUND}(\Arg{x},\Arg{n})
\item Words in \textbf{bold} in the text are defined where they first
  appear and again in the Glossary.
\item A paragraph or a table row marked \DraftSample\ describes a command that
  appears in the menus of this version but is not yet available: choosing
  it displays \texttt{Command not implemented}.  Such a command is listed so
  that the menus in this book match the menus on your screen.
\end{itemize}

\section{Organization of the Initech 123 Reference}
\begin{LotTable}{p{24mm}p{104mm}}
\LotHead{\textbf{Chapter} & \textbf{Contents}}
Chapter 1 & Using Initech 123: the screen, the keys, the menus, entering
  data, ranges, formulas, recalculation and files\\
Chapters 2--12 & The commands of the main menu, one chapter for each, in the
  order the main menu lists them: Add-In, Copy, Data, File, Graph, Move,
  Print, Quit, Range, System and Worksheet\\
Chapter 13 & @Functions, in alphabetical order\\
Appendix A & Worksheet capacity\\
Appendix B & Troubleshooting\\
Glossary & The terms used in this book\LotEnd
\end{LotTable}

Each command chapter opens with a menu tree that shows where its commands
sit in the menus, a table of the commands and what they are for, common uses
of the commands, related reading, and the terms you need to know.  The
command descriptions follow in menu order.

\section{Other Initech 123 Documentation}
\begin{LotTable}{p{50mm}p{78mm}}
\LotHead{\textbf{Manual} & \textbf{Description}}
\doctitle{Setting Up Initech 123, Tutorial, Quick Start and Sample
  Applications} (1051-0000-11) & Four books in one volume: how to start
  Initech 123 on your computer, a series of lessons for new users, a fast
  overview for experienced users, and ready-made worksheets\\
\doctitle{Initech 123 Reference} (1051-0000-10) & This book\\
\doctitle{Initech 123 Quick Reference} (1051-0000-20) & The keys, the
  @functions and the menu trees in a pocket booklet\\
Keyboard templates (1051-0000-21, 1051-0000-22) & The function-key
  assignments, for the 101-key and the 84-key keyboard\LotEnd
\end{LotTable}

\section{After You Complete This Book}
When you have read Chapter~1 you know everything you need to look up any
command in this book.  For practice with a complete worksheet, work through
the \doctitle{Sample Applications} in the Setting Up volume.
